\insimg{forprint/002_folio_01_recto.jpg}
\vskip -10pt

\rightline{\enginline{\small {\it Kuraṇṭakapaddhati}, folio 1\,{\it recto}.}}
\bigskip

\begin{sans}
[ṣaṭkarmāṇi]\\% Section Heading
āsana{}abhyāsena śārīradārḍhye sati ṣaṭkarmāṇi kuryāt\!।

[ādhāraśuddhikriyā]\\ % Subsection Heading: ṣaṭkarma no. 1
vāmadakṣiṇahastatarjanībhyāṃ bhrāmaṇakriyayā śodhayet\!। dine dine bhrāmaṇaṃ\endnote{bhrāmaṇaṃ~] P, brāhmaṇaṃ M.} dvisahasrasaṅkhyāṃ trisahasrasaṅkhyāṃ vā kuryāt\!। śramaparihārārthaṃ tatas $\dagger$tadeva$\dagger$ saṅkhyayā vṛṣapādādikrameṇa āsanāni kuryāt\!। māsatrayādūrdhvaṃ ādhārakambuśuddhaṃ bhavati\!। 

% new paragraph
pathyaṃ bhakṣayet\!। godhūmamudgataṇḍulanīvārakarkaṭikākustumburukṛṣṇamarīcikā ārdrakasaindhava\endnote{ārdraka°~] P, ārdrāka° M.} ityādīni vaidyagranthe prasiddhāni pathyāni\!। dṛḍhe abhyāse jāte tādṛṅniyamo nāsti। bhakṣaṇaṃ tvatitvarayā kartavyam। ānnādīnāṃ rucir na grāhyā vaiguṇyaṃ ca na grāhyam। auṣadhavadannaṃ\endnote{auṣadhavad~] \textit{corr.}, atāṣadhavad M.} bhuñjīta। 

% new paragraph
vighnaparihārārthaṃ vāsudevamantraṃ japet।\endnote{At this point, the scribe has deleted \textit{bhojanottaraṃ oṃkārarahitaṃ japet}.} snānādinā śuddhe sati ekāgrabuddhyā\endnote{ekāgra°~] P, ekrāgra° M.} oṃ namo bhagavate vāsudevāyeti japet। bhojanottaraṃ oṃkārarahitaṃ japet। nidrādau vāsudevavāsuda iti japet। mala-
\end{sans}
\medskip

\begin{center}
\begin{eng}
\suggheading{[Six Cleansing Techniques (ṣaṭkarma)]}
\end{eng}
\end{center}

\begin{multicols}{2}
\begin{eng}
When bodily strength has been achieved by the practice of the postures, [the yogi] should do the
[following] six cleansing techniques.

\headingnoclr{[The Action of Cleansing the Rectum (\textit{ādhāraśuddhikriyā})]}

The [yogi] should cleanse [his rectum] by the action of rotating the fore-fingers of the left and right
hand. Every day, he should do rotations numbering two or three thousand. For alleviating fatigue, he
should then practise postures (\textit{āsana}) according to the sequence beginning with [Pawing] the Leg like a
Bull Pose $\dagger[\ldots]\dagger$. After three months, the rectum and anus are cleansed.

[The yogi] should eat wholesome food. Wheat, mung beans, rice, wild rice, cucumber, coriander, black
pepper, fresh ginger, salt, and the like are well known as wholesome foods in medical literature. When
the practice [of yoga] has become firmly established, there is no such restriction. However, eating
should be done very quickly. One should not dwell on the taste of food and so forth, nor its
imperfections. [The yogi] should eat food as though it were medicine.
For the sake of removing obstacles, [the yogi] should repeat the mantra of Vāsudeva (i.e., Kṛṣṇa). When
he has been purified by bathing and so on, he should repeat with a one-pointed mind, OṂ NAMO
BHAGAVATE VĀSUDEVĀYA. Immediately after food, he should repeat it without the OṂ. Just before sleep,
he should repeat, VĀSUDEVA VĀSUDEVA.
\end{eng}
\end{multicols}
\bigskip

\begin{kan}
\centerline{\suggheadingk{}}
\end{kan}

